\documentclass[../main.tex]{subfiles}
\begin{document}
%==========================================TIÊU ĐỀ TẬP========================================
\begin{table}[h]
    \begin{adjustwidth}{-1cm}{}
        \begin{tabular}{>{\centering\arraybackslash}p{6cm}|>{\raggedright\arraybackslash}p{12.5cm}}
            \multirow{5}{6cm}
            % Cột bên trái: ÉPISODE và số tập
            {\\[-40pt]
                \flaregothic\fontsize{20pt}{20pt}\selectfont \centering
                \color{eptype}HISTOIRE\color{black}\\
                \freshbot\fontsize{64pt}{64pt}\selectfont
                \color{epnum}C\\[6pt]
                \cafeteria\fontsize{20pt}{20pt}\selectfont\color{epnum}(D-C)\color{black}
                \\[18pt]
                \color{black}
            }
             &
            % Dòng thứ nhất: tên tập tiếng Nhật
            {\fotskip\fontsize{18pt}{18pt}\selectfont \ruby{停}{てい}\ruby{電}{でん}の\ruby{夜}{よる}
            }                                                                     \\[6pt]
            % Dòng thứ hai: tên tập tiếng Anh
             & {\cafeteria\fontsize{18pt}{18pt}\selectfont The Night Without Light}                                                                            \\[4pt]
            % Dòng thứ ba: tên tập tiếng Việt
             & {\vnmsans\fontsize{18pt}{18pt}\selectfont Đêm trắng}                                                                                            \\[8pt]
            % Dòng thứ tư: Nhân vật xuất hiện
             & {\caviar{\fontsize{14pt}{14pt}\selectfont Nhân vật: \emp{kuon2} \emp{kaiho2} \emp{ryuuhou2} \emp{suzuno2} \emp{naomi2} \emp{game} \emp{mike1}}}
            \\[8pt]
        \end{tabular}
    \end{adjustwidth}
\end{table}
%=========================================TRUYỆN CHÍNH========================================
\color{black}

\txtbx[2]{
    {\color{vol} \uline{Tóm tắt:}} Cơn mưa rào đầu mùa dữ dội bất ngờ giáng xuống khu phố Anwa, kéo theo một vụ nổ trạm biến áp gây mất điện diện rộng trên toàn khu vực. Giữa căn phòng tầng năm chìm trong bóng tối, gia đình Hikawa cùng nhau quây quần dưới ánh nến lung linh bên nồi súp nghi ngút khói và những câu chuyện cười rộn rã. Cơn mất điện bất chợt gợi lại trong tâm trí Kaiho và Suzuno những ký ức cay đắng về những đêm mưa dột nát, tối tăm ở khu chung cư cũ nơi vũ trụ gốc—nhưng giờ đây, giữa vòng tay chở che của bố mẹ, của anh trai Kuon và các em, bóng tối đã không còn là nỗi sợ, mà trở thành minh chứng cho một mái ấm trọn vẹn mà hai chị em hằng mơ ước.
}

\stpt{-Histoire 1-}
\stcap{Cúp điện bất ngờ}

\pov{Hikawa Kuon}

Khoảng hơn năm giờ chiều một ngày cuối tháng Ba, bầu trời thủ đô Kawachi bỗng chuyển sang một màu xám xịt nặng trĩu. Cơn mưa rào đầu mùa bất ngờ trút xuống như thác đổ, những hạt mưa to bằng hạt ngô quất ràn rạt vào những ô cửa kính lớn của khu nhà ống Hikawa.

Tôi đứng trong gian bếp lớn ở tầng sáu, chuẩn bị nguyên liệu cho bữa cơm tối của cả nhà. Tiếng nước mưa xối xả bên ngoài hòa cùng tiếng sôi lục bục êm tai từ nồi súp củ quả hầm xương trên bếp từ. Bên cạnh tôi, Suzuno đang cẩn thận thái từng lát đậu phụ non mềm mịn xếp ra đĩa, còn bé út Naomi thì ngồi ngoan ngoãn bên chiếc bàn ăn gần đó, cặm cụi tô nốt bức tranh phong cảnh cho bài tập mỹ thuật ở trường.

Ở góc bàn đối diện, Kaiho vừa tan ca xưởng sớm trở về, vẫn còn mặc chiếc áo khoác jeans quen thuộc, đang cặm cụi dùng thước kẻ kiểm tra lại bảng dự trù vật tư kim loại cho bố. Đúng lúc ấy, Ryuuhou từ cầu thang tầng năm lững thững bước lên, tay cầm một quả táo cắn dở, ngó nghiêng nhìn ra ngoài cửa sổ kính: ``Mưa to gió lớn sấm chớp đùng đoàng thế này khiếp thật đấy! May mà hôm nay chị Kaiho về sớm, chứ chậm chân mười phút thôi là ướt như chuột lột ngoài đường rồi.''

Kaiho không ngẩng đầu lên, tay vẫn lia bút chì thoăn thoắt: ``Chị đâu có ngốc như ai kia mà không thèm xem dự báo thời tiết. Thấy mây đen kéo về là chị đã thu dọn xưởng về ngay rồi.''

Vừa dứt lời, từ chiếc đồng hồ thông minh trên cổ tay trái của tôi, một luồng ánh sáng xanh lam dịu nhẹ khẽ chớp lên. Giọng nói máy móc nhưng rành rọt, điềm tĩnh của Game cất lên phá tan tiếng mưa gió: ``Cập nhật dữ liệu từ trạm khí tượng thủy văn Kawachi: Vùng mây đối lưu đang phát triển mạnh mẽ trên toàn bộ lưu vực sông. Dự báo lượng mưa sẽ đạt mức đỉnh điểm trong khoảng chín mươi phút tới, kèm theo giông sét và gió giật cấp sáu. Khuyến cáo mọi cư dân không nên di chuyển ra ngoài nếu không có thiết bị bảo hộ chuyên dụng.''

``Nghe Game nói chưa,'' tôi mỉm cười quay sang các em. ``Tối nay cả nhà đông đủ thế này là tốt rồi. Ngồi quây quần ăn bát súp nóng thì không gì bằng.''

Ryuuhou kéo ghế ngồi cạnh Naomi, vừa nhai táo vừa chăm chú nhìn ngón tay Suzuno thoăn thoắt thái nấm hương. Không gian ngập tràn sự bình yên và ấm áp đặc trưng của một gia đình sum vầy.

Thế nhưng, sự êm đềm ấy chẳng kéo dài được bao lâu.

*XOẸT! PHỤT!!*

Một tia chớp trắng lóa xé toạc bầu trời ngoài cửa sổ, kèm theo tiếng sấm rền vang rung chuyển cả mặt sàn bê tông. Và ngay sau đó, toàn bộ ánh sáng trong căn nhà vụt tắt ngấm trong tích tắc.

Đèn trần tắt ngấm. Màn hình cảm ứng của tủ lạnh tối đen. Chiếc bếp từ đang sôi ùng ục bỗng phát ra một tiếng *tít* dài rồi lịm hẳn. Cả gian bếp tầng sáu rơi vào một bóng tối đặc quánh, chỉ còn lại thứ ánh sáng mờ căm hắt vào từ màn mưa xám xịt ngoài khung cửa.

``\textbf{Á!}'' Naomi giật mình buông rơi chiếc bút màu sáp, khẽ kêu lên một tiếng rồi vội vàng ôm chầm lấy cánh tay của Suzuno.

``Không sao đâu Naomi, có chị ở đây rồi,'' Suzuno vội buông dao, dịu dàng ôm lấy em gái vỗ về.

Tôi đứng yên tại chỗ để mắt làm quen với bóng tối, khẽ cất giọng: ``Game, báo cáo tình trạng đi.''

Màn hình đồng hồ trên tay tôi tự động kích hoạt chế độ tiết kiệm năng lượng, tỏa ra một quầng sáng xanh lục mờ nhạt soi rõ gương mặt của tôi: ``Hệ thống ghi nhận: Toàn bộ nguồn điện lưới cung cấp cho tòa nhà đã bị ngắt do nhảy aptomat tổng tại tầng một. Nguyên nhân sơ bộ: Đường dây trung thế ngoài phố Anwa bị sét đánh lan truyền gây sụt áp đột ngột, kết hợp với tải tiêu thụ thiết bị sưởi và đun nấu trong khu vực tăng đột biến. Nguồn pin dự phòng của tôi còn tám mươi lăm phần trăm, đang ngắt các tiến trình tính toán nền để duy trì kết nối tối thiểu.''

Kaiho đứng dậy, rút chiếc đèn pin mini luôn đeo bên thắt lưng ra bật lên, luồng sáng trắng rọi thẳng vào lối cầu thang: ``Nhảy aptomat tổng à? Để em với anh Kuon xuống tầng một xem sao. Đóng lại cầu dao là xong ngay ấy mà.''

``Em đi với hai người nữa!'' Ryuuhou lập tức xung phong. Con bé lôi chiếc điện thoại trong túi ra, bật đèn flash soi đường rồi lẽo đẽo bước theo sau hai chúng tôi xuống cầu thang.

\ct

Ba anh em men theo lối cầu thang bộ tối om đi qua từng tầng. Không gian tĩnh mịch đến mức chỉ còn nghe thấy tiếng bước chân lộp cộp và tiếng mưa đập xối xả vào tường gạch bên ngoài.

Khi bước qua khúc quanh chiếu nghỉ tầng ba, Ryuuhou bỗng hụt chân một cái, suýt nữa thì va đầu vào thanh lan can sắt nếu tôi không kịp đưa tay đỡ lấy vai con bé.

Tôi nhướng mày nhìn con bé, bật cười trêu: ``Sao thế? Hôm trước đeo kính áp tròng xong bảo nhìn rõ mồn một từng hạt bụi cơ mà?''

Ryuuhou phụng phịu, xoa xoa bên trán, làu bàu tự bào chữa: ``Kính áp tròng nó chỉ giúp điều chỉnh khúc xạ cho mắt nhìn rõ chi tiết thôi, chứ nó có phải kính nhìn đêm hồng ngoại của đặc nhiệm đâu mà đòi soi đường trong bóng tối mù mịt thế này! Trong tối thì mắt mười phần mười cũng vô dụng như mắt cận thị thôi à!''

Kaiho đi phía trước khẽ nhếch môi: ``Thế nên mới bảo bước đi cho cẩn thận, mắt tinh mà chân quáng quàng thì cũng ngã như thường.''

Xuống tới chân cầu thang tầng một, cạnh góc của cửa hàng viễn thông, Kaiho soi đèn pin vào chiếc tủ điện tổng gắn âm tường. Cần gạt aptomat ba pha màu đen đã sập hẳn xuống phía dưới.

Tôi dùng kìm cách điện đẩy cần gạt lên trên đánh *TÁCH* một cái.

Ánh đèn neon ở sảnh tầng một lập tức bừng sáng trở lại. Tiếng quạt thông gió ro ro khởi động. Ryuuhou reo lên một tiếng đầy hy vọng: ``Sáng rồi!''

Thế nhưng, đúng năm giây sau\dots

*TẠCH!!*

Một tia lửa nhỏ lóe lên bên trong hộp kỹ thuật, cần gạt lại nảy ngược trở lại vị trí cũ. Toàn bộ ánh sáng vừa bừng lên lại lịm tắt như chưa từng tồn tại.

Từ chiếc đồng hồ, Game lạnh lùng thông báo: ``Cảnh báo: Trạm biến áp khu vực vẫn chưa ổn định điện áp. Dòng điện dao động bất thường vượt quá ngưỡng an toàn bốn mươi phần trăm. Việc cố tình đóng lại aptomat nhiều lần có thể gây hỏng hóc các vi mạch biến tần của máy móc dưới xưởng và thiết bị gia dụng trên các tầng.''

Kaiho khoanh tay, lắc đầu kết luận: ``Trạm ngoài phố nổ rồi. Tình hình này chắc cả khu Anwa mất điện đến nửa đêm chứ chẳng chơi. Thôi, đừng gạt nữa anh Kuon, kẻo cháy hỏng bo mạch máy CNC của bố dưới xưởng thì khổ.''

Ryuuhou thở dài thườn thượt, ủ rũ tựa lưng vào tường: ``Thế là toi nồi súp tối nay rồi\dots\ Chẳng lẽ tối nay cả nhà phải nhịn đói đi ngủ sớm sao anh chị?''

``Ai bảo nhịn?'' tôi vỗ nhẹ lên vai con bé. ``Nhà mình sống ở khu này bao nhiêu năm, mấy cái chuyện cúp điện này làm sao làm khó được bố con mình. Đi nào, lên tầng năm lấy đồ dự phòng.''

\ct

Mười lăm phút sau, căn phòng sinh hoạt chung ở tầng năm đã được biến đổi thành một không gian hoàn toàn khác lạ.

Biết rằng bếp từ ở tầng sáu không thể hoạt động, tôi đã lấy chiếc bếp ga du lịch mini từ kho dự phòng khẩn cấp ra, đặt ngay ngắn trên chiếc bàn ăn bằng gỗ dày dặn ở tầng năm. Ngọn lửa xanh biếc rì rầm bốc lên, nhẹ nhàng hâm nóng lại nồi súp đang dở dang. Mùi thơm ngọt ngào của nước hầm xương và ngô ngọt lan tỏa khắp căn phòng, xua tan đi cái lạnh lẽo và ẩm ướt của cơn mưa gió bên ngoài.

Trên chiếc bàn trà giữa phòng khách, bốn cây nến sáp thơm mùi quế được thắp sáng, cắm trên những chiếc đĩa sứ trắng muốt.

Ánh lửa nến bập bùng, vàng óng ả hắt lên những bức tường gỗ, tạo nên những vệt bóng lung linh, huyền ảo nhảy múa theo từng cơn gió thoảng qua khe cửa. Không còn ánh đèn huỳnh quang chói gắt, không còn tiếng quạt gió hay tiếng tivi rè rè, cả căn phòng chìm vào một bầu không khí tĩnh mịch, ấm áp và gần gũi đến kỳ lạ.

Cả năm anh chị em chúng tôi quây quần ngồi bệt dưới sàn trải thảm nỉ xung quanh chiếc bàn trà.

Suzuno mở chiếc tủ đồ chơi của Naomi, lôi ra một bộ cờ cá ngựa bằng gỗ cũ kỹ đã phủ một lớp bụi mỏng: ``Tối nay không có tivi hay điện thoại, mấy anh em mình chơi cờ này nhé?''

``Hoan hô!'' Naomi vỗ tay tíu tít, vẻ sợ hãi bóng tối ban nãy đã hoàn toàn biến mất, thay vào đó là sự háo hức rạng rỡ của một đứa trẻ con được cắm trại trong nhà.

Trong lúc hai chị em đang bày những quân cờ gỗ nhỏ xíu ra bàn, Ryuuhou bỗng nảy ra một ý tưởng tinh nghịch. Con bé giơ hai bàn tay lên trước ngọn nến, hướng về phía mảng tường trắng đối diện: ``Này Naomi, nhìn chị làm ảo thuật bóng tay này! Đố biết đây là con gì?''

Em ấy khéo léo đan các ngón tay vào nhau, gập hai ngón cái lại. Trên vách tường, một cái bóng hình đầu một chú chó sói với đôi tai vểnh cao và chiếc mõm dài hiện lên sống động, thi thoảng lại đóng mở miệng như đang sủa.

Naomi thích thú cười khúc khích: ``Con chó sói! Đúng là con chó sói rồi!''

``Đến lượt Naomi này!'' cô bé út hào hứng giơ hai bàn tay nhỏ xíu lên, bắt chéo cổ tay lại rồi vẫy vẫy các ngón tay. Chiếc bóng một chú chim bồ câu trắng đang sải cánh bay lượn rập rờn trên vách tường hiện lên vô cùng mềm mại.

Suzuno mỉm cười dịu dàng, khẽ gập hai đốt ngón tay trỏ và ngón giữa lại, tạo thành chiếc bóng một chú mèo con đang cuộn tròn vẫy đuôi hoàn hảo đến từng chi tiết khiến Naomi ồ lên thán phục.

Thấy cả hai người em đều làm được những hình thù đẹp đẽ, Ryuuhou háo thắng quyết định nâng tầm thử thách: ``Được rồi, xem tuyệt chiêu của chị đây! Đây là một con rồng lửa dũng mãnh trong truyền thuyết!''

Con bé vặn vẹo mười ngón tay đến mức khớp ngón tay kêu răng rắc, hai cánh tay xoắn vặn thành một tư thế kỳ quái hết sức phức tạp.

Kaiho đang ngồi tựa lưng vào mép sofa, liếc nhìn chiếc bóng ngoằn ngoèo, méo mó, thò thụt ba bốn cái chỏm nhọn chẳng ra hình thù gì trên tường, liền bật cười châm chọc: ``Cái củ gừng bị dập nát này là con gì thế Ryuu? Nhìn giống củ khoai tây mọc mầm bị ngập úng thì đúng hơn.''

``Khục\dots!'' Suzuno không nhịn được, vội lấy tay che miệng cười khúc khích.

Ryuuhou mặt đỏ tía tai, hạ tay xuống trừng mắt lườm chị cả: ``Chị Kaiho! Người ta đã bảo đây là con rồng lửa mà! Hai người nhà này chẳng ai có tí tế bào nghệ thuật nào hết á!''

Cả căn phòng rộ lên những tràng cười sảng khoái. Tôi mang ấm trà nóng từ bếp ga mini bước lại bàn, rót ra từng chiếc chén nhỏ cho các em. Hương trà sen thơm ngát ngào ngạt lan tỏa trong không gian mờ ảo của ánh nến.

Thế nhưng, sau tràng cười ấy, tôi nhận thấy ánh mắt của Kaiho và Suzuno bỗng nhiên dừng lại trên ngọn nến bập bùng giữa bàn. Hai chị em khẽ trao nhau một ánh nhìn thật sâu. Một khoảng lặng kỳ lạ, dịu dàng nhưng đong đầy tâm sự bỗng nhiên buông xuống giữa hai người họ.

Suzuno đưa đôi bàn tay nhỏ bé ôm lấy chén trà ấm nóng, ánh mắt nhìn vào làn khói mỏng bốc lên, giọng cô bé nhỏ nhẹ, thì thầm như thể đang trò chuyện với chính quá khứ của mình: ``Tự nhiên\dots\ em lại nhớ đến những ngày tháng ở căn phòng trọ bên thế giới cũ ghê, chị nhỉ\dots''

Kaiho khẽ giật mình, nhưng rồi ánh mắt cô nàng cũng dịu hẳn lại, không còn vẻ gai góc thường ngày. Cô em gái song sinh của tôi khẽ thở dài một hơi nhẹ tênh: ``Ừ\dots\ Cứ mỗi lần trời mưa bão như thế này là cái khu chung cư mini tồi tàn ấy lại cúp điện đầu tiên. Dây điện chắp vá nhằng nhịt, cứ sấm chớp một cái là nhảy cầu dao cả xóm trọ.''

Ryuuhou và Naomi tò mò dừng tay chơi cờ, im lặng lắng nghe. Tôi ngồi bên cạnh, lặng lẽ nhấp một ngụm trà, biết rằng những ký ức sâu kín nhất của hai chị em từ vũ trụ gốc đang chầm chậm ùa về\footnote{{\dawnland HISTOIRE 4}, quyển Khoảng Không.}.

Suzuno khẽ mỉm cười, một nụ cười vừa man mác buồn nhưng lại chan chứa sự ấm áp kiên cường: ``Hồi đó trần nhà dột nát kinh khủng. Mưa to một cái là nước tràn lênh láng qua sàn, nấm mốc mọc khắp nơi. Hai chị em phải gom hết tất cả chậu thau, xô nhựa với mấy chiếc cốc giấy ra đặt khắp phòng để hứng nước. Thế rồi buồn chán quá, hai chị em lại bày trò lấy cốc giấy làm thuyền chơi trò `thủy chiến' ngay trong phòng\dots\ Rồi em trượt chân ngã sặc cả nước, chị Kaiho thì vừa kéo em dậy vừa cười ngặt nghẽo\dots''

``Cười thì cười thế thôi,'' Kaiho tiếp lời, khóe môi khẽ cong lên một nét cười hoài niệm, ``chứ tối đến thì khổ sở trăm bề. Mùa đông lạnh cắt da cắt thịt, hai chị em gom hết mấy tấm chăn cũ rách rưới lại làm một cái tổ trên giường nằm cuộn tròn cho đỡ rét. Điện thoại thì chập chờn, hai đứa chia nhau từng mẩu bánh mì khô khốc với gói mì tôm sống, thắp cây nến rẻ tiền mua ở đầu ngõ để ngồi nghe mấy mẩu chuyện ma nhảm nhí trên mạng\dots''

Giọng em bỗng chùng xuống, mang theo một nỗi niềm cay đắng của những năm tháng một thân một mình gồng gánh nuôi em giữa một xã hội nghiệt ngã: ``Lúc đó, mất điện đối với hai chị em là một cơn ác mộng thật sự. Căn phòng hỏng khóa cửa, chủ trọ thì vô trách nhiệm, xung quanh toàn những kẻ rình rập, cơ hội. Ngồi trong bóng tối mà lúc nào tim cũng đập thình thịch, chỉ sợ có kẻ xấu phá cửa xông vào. Bóng tối ở nơi ấy\dots\ lạnh lẽo, bất an và cô độc đến mức nghẹt thở.''

Căn phòng tầng năm bỗng chìm vào một khoảng lặng sâu lắng. Tiếng mưa rơi lộp độp ngoài khung cửa kính lớn dường như càng làm tôn lên sự tĩnh mịch của căn phòng.

Naomi nghe xong, đôi mắt ngây thơ ngập nước, con bé rướn người sang ôm chặt lấy cánh tay của Kaiho, phụng phịu nói: ``Chị Kaiho với chị Suzuno khổ thân quá\dots\ Nhưng mà bây giờ không sợ nữa đâu ạ! Ở đây có Anh Kuon, có chị Ryuuhou, có bố Ryujin, có Naomi bảo vệ các chị rồi!''

Lời nói ngây ngô của cô bé út khiến lớp băng giá cuối cùng trong lòng Kaiho tan biến hoàn toàn. Cô nàng cúi xuống xoa đầu Naomi, rồi đưa bàn tay thô ráp của một người thợ cơ khí khẽ nắm lấy bàn tay mềm mại của Suzuno bên cạnh.

Suzuno ngước nhìn chị gái, mỉm cười thật tươi, đôi mắt lấp lánh phản chiếu ánh lửa nến: ``Vâng ạ. Ở đây ấm áp lắm. Cũng là bóng tối, cũng là thắp nến trong đêm mưa bão\dots\ nhưng ở nơi này, em không còn cảm thấy sợ hãi hay cô đơn nữa. Vì em biết, chị em mình đã thực sự có một gia đình trọn vẹn rồi.''

Tôi nhìn hai chị em họ, trong lòng dâng lên một cảm giác xúc động nghẹn ngào. Tôi vươn tay đặt nhẹ lên vai Kaiho và Suzuno, khẽ nói: ``Dù quá khứ có vỡ vụn đến đâu, thì từ nay về sau, ngôi nhà này sẽ luôn là nơi che chở cho hai đứa. Bão giông hay mất điện ngoài kia sẽ không bao giờ có thể làm tổn thương gia đình mình được nữa đâu.''

Kaiho liếc nhìn tôi, ánh mắt chan chứa sự tin cậy và biết ơn sâu sắc mà hiếm khi cô nàng chịu bộc lộ bằng lời nói.

\ct

Để xua tan bầu không khí có phần trầm lắng, Ryuuhou bỗng hắng giọng, hai tay chống cằm nhìn ngọn nến bập bùng: ``Nãy chị Suzuno nhắc chuyện ma làm em lại tò mò ghê\dots\ Anh Kuon! Nhân dịp cúp điện lãng mạn thế này, anh kể một câu chuyện ma cho cả nhà nghe đi!''

Suzuno đang tựa đầu vào vai Kaiho, nghe vậy liền lắc đầu khẽ can: ``Thôi nào Ryuu, Naomi sắp buồn ngủ rồi, em cứ thích nghe mấy chuyện rùng rợn dọa ma dọa quỷ.''

``Ai bảo chuyện ma là rùng rợn?'' tôi mỉm cười, nhấp một ngụm trà sen. ``Anh chỉ biết một câu chuyện cổ xưa về những ngôi nhà cổ ở vùng đồng bằng Bắc Bộ thôi.''

Ryuuhou lập tức ngồi thẳng lưng dậy, hai mắt sáng rực: ``Kể đi anh, em hứa là không sợ đâu!''

Tôi im lặng mất chừng mười giây, để căn phòng hoàn toàn chìm vào tiếng mưa gió xào xạc bên ngoài. Rồi bằng một chất giọng trầm ấm, chậm rãi và ma mị hiếm thấy, tôi bắt đầu kể.

``Hồi anh còn nhỏ, bố Ryujin từng kể rằng, những ngôi nhà ống cổ thường được dựng bằng những thân gỗ lim đại thụ ngàn năm tuổi. Người xưa tin rằng gỗ lim sau khi đốn hạ vẫn giữ lại linh hồn của rừng sâu. Vào những đêm mưa bão sấm sét mất điện, khí lạnh từ lòng đất bốc lên làm thớ gỗ co thắt, đánh thức linh hồn ấy cựa mình thức giấc\dots\ Khi cả ngôi nhà chìm vào giấc ngủ say, trên căn gác xép cao nhất sẽ bắt đầu vang lên những âm thanh kỳ lạ\dots\ Tiếng bước chân trần kéo lê lẹt xẹt, rồi sau đó\dots\ là tiếng gõ cửa buốt lạnh từng nhịp vang lên từ khoảng không vô định\dots''

Đúng lúc câu nói của tôi vừa dứt\dots

*CỐC! CỐC! CỐC!*

Một chuỗi âm thanh khô khốc, đanh gọn vang lên rõ mồn một ngay trên trần nhà, từ phía hành lang tầng tám dội thẳng xuống.

Không gian xung quanh như đóng băng lại trong tích tắc.

Ryuuhou giật nảy bắn người, nhảy dựng khỏi sàn nhà, mặt cắt không còn một giọt máu, hai tay ôm chặt lấy đầu hét toáng lên: ``\textbf{Áaaaa! Cái gì thế kia?! Có ma thật kìa Onii ơi!!}''

Ngay cả Kaiho và Suzuno cũng khẽ giật mình thót lại, đưa mắt nhìn chăm chú lên trần nhà tối om.

Đúng lúc Ryuuhou đang run rẩy nép sát vào lưng tôi thì từ trên chiếc đồng hồ thông minh, giọng nói máy móc của Game lại bình thản vang lên như một gáo nước lạnh dội thẳng vào cơn hoảng loạn: ``Phân tích âm học hoàn tất: Nguồn âm thanh phát ra từ hệ thống đường ống dẫn nước kim loại gắn dọc hộp kỹ thuật tầng tám. Do sự chênh lệch nhiệt độ đột ngột giữa nước mưa lạnh bên ngoài và không khí ấm bên trong hộp gen, các khớp nối ren kim loại đã co ngót đột ngột tạo ra hiện tượng búa nước cơ học.''

Sự thật khoa học rành rọt của Game khiến bầu không khí ma quái bỗng chốc vỡ tan tành.

Ryuuhou đứng hình mất vài giây. Nhìn thấy nụ cười nhịn không nổi đang nở toe toét trên gương mặt tôi và ánh mắt trêu chọc của Kaiho, cô nhóc mới vỡ lẽ mình vừa bị trúng kế lừa ngoạn mục.

Con bé tức tối túm lấy chiếc gối tựa ném thẳng vào ngực tôi, vừa giậm chân vừa mếu máo cằn nhằn: ``Anh Kuon!! Tại sao anh lại căn đúng cái lúc ống nước gõ cành cạch để kể câu chuyện ma quái ấy hả?! Anh làm em suýt nữa thì rớt tim ra ngoài rồi đây này!!''

Tôi ôm chiếc gối bông bật cười thành tiếng, còn Kaiho thì lắc đầu châm chọc: ``Đấy, đã nhát gan mà cứ thích ra vẻ anh hùng. Nghe ống nước gõ mà tưởng ma nhập trần nhà!''

Căn phòng tầng năm lại tràn ngập những tiếng cười khúc khích xua tan đi màn đêm tăm tối.

\ct

Khoảng hơn chín giờ tối, cơn mưa rào ngoài trời đã hoàn toàn dứt hạt.

Bé út Naomi đã say giấc nồng từ lâu, được Suzuno bế vào chiếc giường đệm ấm áp ở góc phòng tầng năm, đắp chiếc chăn bông mỏng ngang ngực. Suzuno sau một hồi ngồi đọc sách dưới ánh nến cũng bắt đầu díp mắt lại, dựa đầu vào thành giường ngủ thiếp đi. Ryuuhou thì nằm cuộn tròn trên chiếc ghế sofa dài, hơi thở đều đặn và êm ái.

Tôi và Kaiho thu dọn những chiếc chén trà trên bàn. Nhìn ngọn nến duy nhất còn cháy le lắt, tôi khẽ nói với Kaiho: ``Anh lên kiểm tra tầng tám một chút nhé.''

Kaiho gật đầu: ``Đi cẩn thận, cầu thang trơn đấy.''

Tôi cầm ngọn nến cắm trên chiếc đĩa sứ nhỏ, nhẹ nhàng bước từng bước êm ái theo cầu thang bộ đi lên các tầng trên.

Tầng sáu - gian bếp lớn và phòng thờ tổ tiên - tất cả các cánh cửa chớp đều đã được đóng chặt, không có hạt mưa nào lọt vào.

Tầng bảy - căn hộ của Kson - cánh cửa gỗ đóng im lìm, bên trong hoàn toàn yên ắng.

Tôi tiếp tục bước lên những bậc cầu thang cuối cùng dẫn lên tầng tám.

Khi bước chân tôi dừng lại ở chiếu nghỉ tầng cao nhất, ánh mắt tôi bất giác nhìn về phía cánh cửa phòng của Mikeneko.

Dưới khe cửa hẹp sát mặt sàn bê tông, một vệt sáng màu vàng cam ấm áp, run rẩy đang hắt ra ngoài hành lang tối tăm.

Ánh nến.

Trong căn phòng ấy, Mikeneko vẫn chưa ngủ. Cô nàng cũng đang ngồi một mình trong bóng tối, chỉ với một ngọn nến nhỏ leo lắt bầu bạn giữa màn đêm mất điện tịch mịch.

Tôi đứng lặng yên ở hành lang một lúc lâu. Tiếng gió đêm sau mưa thổi qua ô thoáng hành lang mang theo hơi ẩm ngai ngái dịu mát.

Tôi không bước tới gõ cửa. Tôi biết tính cách khép kín và nhạy cảm của cô ấy; một tiếng gõ cửa bất ngờ vào đêm muộn như thế này chỉ khiến cô nàng thêm hoảng sợ và co cụm lại vào vỏ ốc của mình.

Tôi quay người, nhẹ nhàng bước xuống bếp tầng sáu. Dưới ánh lửa xanh liu riu của chiếc bếp ga du lịch, tôi pha một tách trà hoa cúc mật ong nóng hổi vào chiếc cốc sứ có nắp đậy giữ nhiệt.

Cầm cốc trà nóng bốc khói nghi ngút trên tay, tôi quay trở lại tầng tám.

Tôi bước thật khẽ, cúi người đặt chiếc cốc trà ấm áp ngay ngắn xuống mặt sàn trước cửa phòng Mikeneko, ngay cạnh vị trí vệt sáng ánh nến hắt ra. Tách trà tỏa ra hương thơm dịu ngọt của mật ong và hoa cúc thanh tao, như một lời nhắn nhủ ấm áp và thầm lặng giữa đêm đen lạnh giá.

Xong xuôi, tôi không nán lại thêm một giây nào. Tôi cầm ngọn nến của mình, nhẹ nhàng xoay người bước xuống cầu thang, trả lại cho tầng tám sự tĩnh lặng vốn có.

\vt

Gần mười một giờ đêm, khi tôi và Kaiho đang ngồi trò chuyện khe khẽ bên bàn trà dưới ánh nến tàn, thì từ trên trần nhà bỗng vang lên một tiếng *tách* quen thuộc.

*OÒNG!!*

Hệ thống điện lưới khu vực rốt cuộc cũng được khôi phục.

Toàn bộ bóng đèn hành lang và đèn trần trong phòng khách đồng loạt bừng sáng trở lại. Tiếng ro ro êm ái của máy lọc không khí và tiếng chuông báo hoàn thành chu trình của tủ lạnh vang lên rộn rã.

Từ cổ tay tôi, Game cất giọng thông báo: ``Trạm biến áp trung thế đã hoàn tất sửa chữa. Điện áp ổn định ở mức hai trăm hai mươi vôn. Tải điện an toàn.''

Tôi và Kaiho vội vàng đứng dậy, tắt bớt những công tắc đèn trần công suất lớn, chỉ để lại dải đèn hắt tường màu vàng nhạt dịu mắt để không làm kinh động đến giấc ngủ của các em.

Naomi vẫn đang ôm chặt chú gấu bông ngủ say sưa, khóe môi khẽ chu ra như đang mơ thấy một giấc mơ ngọt ngào. Suzuno khẽ trở mình, kéo chiếc chăn mỏng đắp kín vai cho em gái. Còn Ryuuhou thì gác chân chữ ngũ lên thành ghế sofa, miệng lẩm bẩm vài câu vô nghĩa về mấy con rồng lửa và que kem đậu xanh.

Kaiho đứng bên cạnh tôi, cúi đầu thổi tắt ngọn nến cuối cùng trên bàn trà. Một làn khói trắng mỏng manh thoang thoảng mùi quế bay lên rồi tan biến vào không gian.

Cô nàng nhìn ngắm các em đang say ngủ, rồi khẽ quay sang tôi, khóe môi nở một nụ cười hiếm hoi, chân thành và ấm áp: ``Cảm ơn anh vì bữa tối dưới ánh nến nhé, Kuon.''

Tôi mỉm cười, khẽ vỗ nhẹ vai người em gái song sinh: ``Đêm nay là một đêm rất đẹp mà, đúng không?''

Cả hai chúng tôi cùng bước lại gần cửa sổ, kéo tấm rèm che bớt ánh sáng đèn đường rọi vào phòng, khép lại một đêm trắng đầy ắp những kỷ niệm thiêng liêng và dịu dàng của gia đình Hikawa.

\omk

Sáng ngày hôm sau, bầu trời Kawachi lại quang đãng, trong xanh không một gợn mây.

Khoảng hơn bảy giờ sáng, trước cửa phòng Mikeneko ở tầng tám, chiếc cốc sứ đựng trà hoa cúc mật ong đêm qua đã được uống cạn sạch, được rửa sạch sẽ và đặt ngay ngắn trên một chiếc khăn giấy khô, bên cạnh là một mẩu giấy nhỏ vẽ hình một chú mèo con đang cười tít mắt thay cho lời cảm ơn.

Dưới phòng sinh hoạt tầng năm, Ryuuhou vừa thức dậy, đầu tóc bù xù bước ra bàn ăn, vẫn chưa hết ấm ức vụ tối qua liền chống nạnh càu nhàu: ``Sáng nay em phải đi kiểm tra lại cái ống nước tầng tám mới được! Ai đời đang lúc nghe chuyện ma rùng rợn thì nó lại gõ đúng ba tiếng như có người gõ cửa thật, làm em mất hết cả thể diện trước mặt bé Naomi!''

Kaiho đang ngồi nhâm nhi tách cà phê nóng, liếc nhìn cô em gái ngổ ngáo, rồi bật cười trêu: ``Anh Kuon lúc nào chẳng biết cách tính thời điểm chuẩn xác như một cái đồng hồ Thụy Sĩ. Nhưng mà nói cho em biết nhé: Bóng tối chỉ làm cho những đứa nhát cáy như em sợ hãi thôi. Còn với những người từng trải qua đêm tối thực sự, thì ánh nến tối qua chính là thứ ánh sáng ấm áp và đẹp đẽ nhất trần đời đấy.''

Suzuno bưng đĩa trứng ốp la nóng hổi đặt lên bàn, khẽ mỉm cười đồng tình với chị cả.

Ryuuhou chớp chớp mắt, dường như hiểu ra ý nghĩa sâu xa trong câu nói của chị, con bé không cãi lại nữa mà chỉ cười toe toét, nhảy tót vào bàn ăn chén một miếng bánh mì giòn rụm, đón chào một ngày mới tràn ngập ánh nắng ban mai.

\end{document}
